\begin{frame}{``예상대로 졌다''는 예상된 결과}
  \begin{itemize}
    \item digits 같은 \textbf{작은 정형} 데이터에서는 \textbf{신경망이
      GBDT를 (소폭) 이기지 못하는 것이 \textit{예상되는} 결과}
    \item Grinsztajn, Oyallon \& Varoquaux (2022), ``Why do tree-based
      models still outperform deep learning on typical tabular data?'':
      \textbf{저차원 정형 데이터에서는 GBDT가 신경망을 이긴다}
    \item 신경망의 우위는 \textbf{세 조건} 중 하나에서 온다 —
      (a) \textbf{차원이 높은} 데이터(텍스트를 원시 토큰으로 보면
      $10^4$$\sim$$10^5$차원), (b) \textbf{데이터가 아주 큰} 경우,
      (c) 특징이 \textbf{공유되는} 국소 구조(이미지)
  \end{itemize}
  \vskip0.5em
  \begin{block}{이 프로젝트의 태도}
    digits는 64차원(저차원) $\Rightarrow$ GBDT가 앞서는 쪽이
    \textit{예상}이고, 실험은 정확히 그렇게 나온다 (64차원에서
    GBDT 0.944 $\succ$ CNN 0.939, 격차 0.006).
    \kq{``왜 DL''의 정직한 답: ``신경망이 좋다''가 아니라
    ``\textbf{언제}(어떤 차원·구조에서) 신경망이 좋은지''를
    실증한다.}
  \end{block}
\end{frame}
